\scanpage{13}
% archive sheet number: 12
\mspage{- IV-\handwritten{\unsure{963|96}} -}

\uline{briques séparables de} $\handwritten{\kappa}(\underline p)$ (ce qui entraîne que les fibres du \unsure{morphisme} canon\handwritten{i}que $\operatorname{Spec}({}^{\mathrm{hs}}A)\handwritten{\to}\operatorname{Spec}(A)$ sont \uline{géométriquement régulières}) .

\uline{Corollaire} (18.8.\deleted{9}\handwritten{12}) .- \uline{Soit} $A$ \uline{un anneau local noethérien . Pour que} $ {}^{\mathrm{hs}}A$ \uline{possède une des propriétés suivantes}\deleted{\unsure{,}} :

a) \uline{être un anneau de Cohen-Macaulay} ;

b) \uline{vérifier la condition} $(S_n)$ ;

c) \uline{être régulier} ;

d) \uline{vérifier la condition} $(R_n)$ ; \margin{\handwritten{\tikz{\draw (0,0) .. controls (.05,-.4) and (.07,-.5) .. (.18,-.8); \draw (-.65,-.27) .. controls (-.65,-1) and (-.31,-.81) .. (.18,-.8); \draw (-.65,-1.25) .. controls (-.54,-.90) and (-.22,-.78) .. (.18,-.8);} \unsure{11|''}\,$\vert$}}

\uline{il faut et il suffit que $A$ possède cette même propriété} .

\uline{Corollaire} (18.8.1\deleted{0}\handwritten{3}) .- \uline{Soit} $A$ \uline{un anneau semi-local noethérien . Pour qu'un idéal premier} $\underline p$ \uline{de} $ {}^{\mathrm{hs}}A$ \uline{appartienne à} $\operatorname{Ass}({}^{\mathrm{hs}}A)$ , \uline{il faut et il suffit que} $\underline p\handwritten{\in} A$ \uline{appartienne à} $\operatorname{Ass}(A)$ .

Les démonstrations sont les mêmes que celles de (18.6.9),(18.6.10) et (18.6.11) respectivement .

\uline{Proposition} (18.8.1\deleted{1}\handwritten{4}) .- \uline{Soit} $A$ \uline{un anneau local ; les propriétés suivantes sont équivalentes} :

a) $A$ \uline{est géométriquement unibranche} (resp. \uline{réduit et géométriquement unibranche}) .

b) \uline{Pour toute} $A$-\uline{algèbre locale essentiellement étale} $A'$ , $\operatorname{Spec}(A')$ \uline{est irréductible} (resp. \uline{intègre}) .

c) $\operatorname{Spec}({}^{\mathrm{hs}}A)$ \uline{est irréductible} (resp. \uline{intègre}) .

On se ramène comme dans (18.6.1\handwritten{2}) au cas où $A$ est réduit \handwritten{[}, \uwave{en utilisant (18.8.7.1) ,}\handwritten{]}\margin{\handwritten{\unsure{!|?}}} et l'équivalence de b) et c) se prouve comme dans (18.6.1\handwritten{2}) ; i\handwritten{l} en est de même du fait que a) entraîne c) , en tenant compte de (18.8.7.1) et de la définition d'un anneau local intègre géométriquement unibranche . Enfi\handwritten{n} pour démontrer que c) entraîne a) , gardons les mêmes notations que dans (18.6.1\handwritten{2}) ; il s'agit encore de voir que $ {}^{\mathrm{hs}}B$ est intègre . Mais , par (18.8.7) , $ {}^{\mathrm{hs}}B$ est un anneau local de l'anneau semi-local $B\handwritten{\otimes}_A({}^{\mathrm{hs}}A)$ , qui s'identifie encore par platitude à un sous-anneau du corps $L$ , donc est intègre ; il en est par suite de même de $ {}^{\mathrm{hs}}B$, \handwritten{\unsure{/}}ce qui achève la démonstration .

\deleted{\unsure{Proposition}}\inserted{Corollaire} (18.8.1\deleted{2}\handwritten{5}) .- \uline{Soit} $A$ \uline{un anneau \inserted{local} hensélien} (resp\handwritten{.} \uline{strictement local}\handwritten{)}
